\chapter{Heat Kernel 2 (in curved space)}
 The heat kernel function is defined as the solution of the following differential equation
\begin{equation}\label{heat equation}
        \frac{\partial\,\mathcal{G}(s;x,x')}{\partial s} + \mathcal{O}_x\mathcal{G}(s;x,x')=0
\end{equation}
with initial condition
\begin{equation}
    \mathcal{G}(0;x,x')=\delta^{(d)}(x,x')
\end{equation}
where $\delta^{(d)}(x,x')$ is the biscalar $\delta$-function that generalizes the usual flat space Dirac delta, and $x,x'$ are (coordinates of) points on an euclidean (i.e. riemannian) manifold of dimension $d$. To extend the results to the usual spacetime it is necessary to assume an analytic continuation of any Minkowski metric to one of euclidean signature.
Restricting our attention to simple scalar fields, we can consider an operator $\mathcal{O}_x$ of Laplace-type
\begin{equation}\label{Laplace-type operator}
    \mathcal{O} = - g^{\mu\nu}\nabla_\mu \partial_\nu + E
\end{equation}
in which $E=E(x)$ is a local endomorphism acting multiplicatively on the scalar field's bundle. If we solve the diffusion equation (\ref{heat equation}) implicitly
\begin{equation}
\mathcal{G}(s;x,x')= \bra{x'}e^{-s\mathcal{O}}\ket{x},
\end{equation}
we can see that heat kernel function is related to the Green function $G$, which is formally defined by
\begin{equation}
    \mathcal{O}_x G(x,x')= \delta^{(d)}(x,x').
\end{equation}
Thus,
\begin{equation}
    G(x,x')= \int_0^\infty ds\,\mathcal{G}(s;x,x').
\end{equation}
The heat kernel function has an asymptotic expansion for $s \xrightarrow{} 0^+$. Following DeWitt \cite{DeWitt:1964mxt}, it has the form
\begin{equation}\label{Seely-DeWitt expansion}
    \mathcal{G}(x,x') = \frac{\Delta(x,x')^{1/2}}{(4 \pi s)^{d/2}} e ^{-\frac{\sigma(x,x')}{2s}}\sum_{k \geq 0} a_k (x,x') s^k.
\end{equation}
In Eq.(\ref{Seely-DeWitt expansion}) several bitensors are introduced, the most fundamental is $\sigma(x,x')$ called \emph{geodetic interval} or \emph{Synge-De Witt's world function}. It is defined as half of the square of the geodesic distance between $x$ and $x'$. The bitensor $\Delta(x,x')$ is known as \emph{van Vleck determinant} and is related to the world function and the determinant metric by
\begin{equation}
    \Delta(x,x') = \frac{1}{g(x)^{1/2}\,g(x')^{1/2}}\;\text{det}\,\left(-\frac{\partial^2}{\partial x^\alpha \partial x^{'\beta}}\sigma(x,x')\right).
\end{equation}
Together, $\sigma$ and $\Delta$ ensure that the leading term of the Seely-De Witt parametrization covariantly generalizes the solution of the heat equation in flat space, where $\mathcal{O} \sim -\partial^2.$ Finally, the bitensors $a_k(x,x')$ are the coefficients of the asymptotic expansion and contain the geometrical information of the operator $\mathcal{O}$. The ultraviolet properties are local in renormalizable theories and for the case of the heat kernel locality correspond to $x \sim x'$ and it is captured by the coincidence limit in which $x' \xrightarrow{} x$. Given any bitensor $B(x,x')$, its coincidence limit is defined 
\begin{equation}
    [B] \coloneqq \lim_{x' \xrightarrow{} x} B(x,x').
\end{equation}
One important note is that covariant derivatives do not generally commute with the coincidence limit, so 
\begin{equation}
    \nabla  [B] \neq [\nabla B].
\end{equation}
The coincidence limits of the bitensors $\sigma(x,x')$ and $\Delta(x,x')$ and their derivatives can be obtained by repeated differentiation of the \emph{crucial relations}, obtained in \cite{DeWitt:1964mxt} by geometrical observations for geodesics, they read
\begin{equation}\label{crucial relations}
    \sigma_\mu \sigma^\mu = 2 \sigma,\;\;\;\;\;\; \Delta^{1/2} \sigma_\mu^{\;\,\mu} + 2 \sigma ^\mu \nabla_\mu \Delta^{1/2} = d \Delta^{1/2},
\end{equation}
in which we suppressed the bitensor coordinates and we introduced the notation in which subscripts of $\sigma$ indicate covariant derivative, i.e. $\sigma_{\mu_{1}\dots\mu_{n}} \coloneqq \nabla_{\mu_{n}}\dots\nabla_{\mu_{1}} \sigma$. We can start with the first of Eq.(\ref{crucial relations}) obtaining
\begin{align}
\sigma_{\nu} &= \sigma_{\mu \nu} \sigma^{\mu}\label{diff sigma 0}&&\\
\sigma_{\nu\rho} &= \sigma_{\mu\nu\rho}\sigma^\mu + \sigma_{\mu\nu} \sigma^{\mu}_{\;\,\rho}\label{diff sigma 1}&&\\
\sigma_{\nu\rho\sigma} &= \sigma_{\mu\nu\rho\sigma} \sigma^\mu + \sigma_{\mu\nu\rho} \sigma^\mu_{\,\;\sigma} + \sigma_{\mu\nu\sigma} \sigma^\mu_{\,\;\rho} +\sigma_{\mu\nu} \sigma^\mu_{\,\;\rho\sigma}\label{diff sigma 2}&&\\
\sigma_{\nu\rho\sigma\tau} &= \sigma_{\mu\nu\rho\sigma\tau} \sigma^\mu + \sigma_{\mu\nu\rho\sigma} \sigma^\mu_{\,\;\tau} + \sigma_{\mu\nu\rho\tau}\sigma^\mu_{\,\;\sigma} + \sigma_{\mu\nu\rho}\sigma^\mu_{\,\;\sigma\tau} + \sigma_{\mu\nu\sigma\tau}\sigma^\mu_{\,\;\rho}\label{diff sigma 3} \\
&\quad + \sigma_{\mu\nu\sigma}\sigma^\mu_{\,\;\rho\tau} +\sigma_{\mu\nu\tau} \sigma^\mu_{\,\;\rho\sigma} +\sigma_{\mu\nu}\sigma^\mu_{\,\;\rho\sigma\tau}\notag&&\\
&\;\;\vdots \notag
\end{align}
Now we can compute coincidence limits considering that we already know
\begin{equation}\label{c.l.sigma 1}
    [\sigma] \coloneqq \lim_{x' \xrightarrow{} x} \sigma(x,x') = 0.
\end{equation}
Using (\ref{c.l.sigma 1}) and the first of (\ref{crucial relations}) we get also
\begin{equation}\label{c.l.sigma 2}
    [\sigma_\mu] = 0.
\end{equation}
We notice that using (\ref{c.l.sigma 2}), the coincidence limit of (\ref{diff sigma 0}) is a trivial identity, while for (\ref{diff sigma 1}) we have
\begin{equation}\label{c.l.sigma 3}
    [\sigma_{\mu\nu}] = g_{\mu\nu}.
\end{equation}
To compute coincidence limits for (\ref{diff sigma 2}) and (\ref{diff sigma 3}), we need to use some commutation laws for covariant differentiation together with (\ref{c.l.sigma 1}),(\ref{c.l.sigma 2}) and (\ref{c.l.sigma 3}). We have
\begin{equation}
    [\sigma_{\nu\rho\sigma}]= \lim_{x \xrightarrow{} x'} (\sigma_{\nu\rho\sigma} + \sigma_{\rho\nu\sigma} + \sigma_{\sigma\nu\rho})
\end{equation}
and we now subtract $3[\sigma_{\nu\rho\sigma}]$ and use the Ricci identity
\begin{equation}
    \begin{split}
        -2[\sigma_{\nu\rho\sigma}] &= [\sigma_{\rho\nu\sigma}] - [\sigma_{\nu\rho\sigma}] + [\sigma_{\sigma\nu\rho}] - [\sigma_{\nu\rho\sigma}]\\
        &= ([\sigma_{\rho\nu\sigma}] - [\sigma_{\nu\rho\sigma}]) + ([\sigma_{\sigma\nu\rho}] -[\sigma_{\nu\sigma\rho}]) + ([\sigma_{\nu\sigma\rho}] - [\sigma_{\nu\rho\sigma}])\\
        &= [\nabla_{\sigma}[\nabla_\nu,\nabla_\rho]\sigma] + [\nabla_\rho[\nabla_\nu,\nabla_\sigma]\sigma] + [R_{\rho\sigma\nu}^{\;\;\;\;\;\mu} \sigma_\mu] = 0,
    \end{split}
\end{equation}
where in the last line the first two terms disappear because we assume a torsion-free theory, and the last term because of (\ref{c.l.sigma 2}), so 
\begin{equation}\label{c.l.sigma 4}
    [\sigma_{\nu\rho\sigma}] = 0.
\end{equation}
A consequence of considering a torsion-free theory is reflected by the symmetry for the first two indexes closest to $\sigma$, namely $[\sigma_{\mu\nu\dots} ] = [\sigma_{\nu\mu\dots}]$.
For the four-derivative term, using (\ref{c.l.sigma 4}), we have
\begin{equation}
    [\sigma_{\nu\rho\sigma\tau}] = [\sigma_{\tau\nu\rho\sigma}] + [\sigma_{\sigma\nu\rho\tau}] + [\sigma_{\rho\nu\sigma\tau}] + [\sigma_{\nu\rho\sigma\tau}].
\end{equation}
Now we can subtract $4[\sigma_{\nu\rho\sigma\tau}]$ and use the relation
\begin{equation}
    [\nabla_\mu,\nabla_\nu]T_{\rho\sigma} = R_{\mu\nu\rho}^{\;\;\;\;\lambda}T_{\lambda\sigma} + R_{\mu\nu\sigma}^{\;\;\;\;\lambda}T_{\rho\lambda},
\end{equation}
which, in the case of $T_{\mu\nu} = \sigma_{\mu\nu}$, gives simply
\begin{equation}
    [\nabla_\mu,\nabla_\mu] \sigma_{\rho\sigma}=0.
\end{equation}
We get
\begin{equation}
    \begin{split}
        -3[\sigma_{\nu\rho\sigma\tau}] &= [\sigma_{\tau\nu\rho\sigma}] - [\sigma_{\nu\rho\sigma\tau}] + [\sigma_{\sigma\nu\rho\tau}] - [\sigma_{\nu\rho\sigma\tau}] + [\sigma_{\rho\nu\sigma\tau}] - [\sigma_{\nu\rho\sigma\tau}]\\
        &= [\sigma_{\tau\nu\rho\sigma}] - [\sigma_{\nu\rho\sigma\tau}] + [\sigma_{\nu\sigma\rho\tau}] - [\sigma_{\nu\rho\sigma\tau}]\\
        &= ([\sigma_{\tau\nu\rho\sigma}] - [\sigma_{\nu\rho\tau\sigma}]) + ([\sigma_{\nu\rho\tau\sigma}] +[\sigma_{\nu\rho\sigma\tau}]) + ([\sigma_{\nu\sigma\rho\tau}] - [\sigma_{\nu\rho\sigma\tau}])\\
        &= [\nabla_\sigma(R_{\rho\tau\nu}^{\;\;\;\;\mu}\sigma_{\mu})] + [\nabla_\tau(R_{\rho\sigma\nu}^{\;\;\;\;\mu}\sigma_{\mu})] +[[\nabla_\sigma.\nabla_\tau]\sigma_{\rho\nu}]\\
        &= R_{\rho\tau\nu\sigma} + R_{\rho\sigma\nu\tau}.
    \end{split}
\end{equation}
Using symmetries of the Riemann tensor we get
\begin{equation}\label{c.l.sigma 5}
    [\sigma_{\nu\rho\sigma\tau}] = \frac{1}{3}(R_{\nu\sigma\tau\rho} + R_{\nu\tau\sigma\rho}).
\end{equation}
Other relations can be obtained taking contractions in (\ref{diff sigma 3}) and performing further differantiation, as
\begin{equation}
    [\sigma_{\nu\,\;\sigma\,\;\rho}^{\,\;\nu\,\;\sigma}] = \nabla_\rho R
\end{equation}
and
\begin{equation}
    [\sigma_{\mu\,\;\nu\,\;\sigma}^{\,\;\mu\,\;\nu\,\;\sigma}]=\frac{8}{5}\nabla^\mu\nabla_\mu R + \frac{4}{15} R_{\mu\nu} R^{\mu\nu} - \frac{4}{15}R_{\mu\nu\rho\sigma}R^{\mu\nu\rho\sigma}
\end{equation}
We can now analyze $\Delta$ in the same fashion, differentiating the second of (\ref{crucial relations}), we get
\begin{align}
    d\nabla_\nu \Delta^{1/2} &= \nabla_\nu \Delta^{1/2} \sigma_\mu^{\;\mu} + \Delta^{1/2} \sigma_{\mu\;\nu}^{\;\mu} + 2 \sigma^\mu_{\;\nu} \nabla_\mu \Delta^{1/2} + 2\sigma^\mu \nabla_\nu \nabla_\mu \Delta^{1/2}\label{diff delta 0}\\
    d\nabla_\rho \nabla_\nu \Delta^{1/2} &= \nabla_\rho\nabla_\nu \Delta^{1/2} \sigma_\mu^{\;\mu} + \nabla_\nu\Delta^{1/2} \sigma_{\mu\;\rho}^{\;\mu} + \nabla_\rho \Delta^{1/2} \sigma_{\mu\;\nu}^{\;\mu} +\Delta^{1/2} \sigma_{\mu\;\nu\rho}^{\;\mu}\label{diff delta 1}\\
    &+ 2 \sigma^\mu_{\;\nu\rho} \nabla_\mu \Delta^{1/2} +2 \sigma^\mu_{\;\nu} \nabla_\rho \nabla_\mu \Delta^{1/2} + 2\sigma^\mu_{\;\rho} \nabla_\nu \nabla_\mu \Delta^{1/2} + 2\sigma^\mu \nabla_\rho\nabla_\nu \nabla_\mu \Delta^{1/2}. \notag
\end{align}
Before computing coincidence limits we notice that by definition we have
\begin{equation}\label{c.l.delta 0}
    [\Delta] = 1.
\end{equation}
In Eq.(\ref{diff delta 0}), using coincidence limits of $\sigma$, we get trivially
\begin{equation}\label{c.l.delta 1}
    [\nabla_\nu \Delta^{1/2} ] = 0.
\end{equation}
In Eq.(\ref{diff delta 1}), in particularly using (\ref{c.l.sigma 5}) we get
\begin{equation}\label{c.l.delta 2}
    [\nabla_\rho \nabla_\nu \Delta^{1/2}] = \frac{1}{6} R_{\rho\nu},
\end{equation}
and similarly
\begin{equation}
    [\nabla_\rho \nabla_\nu \Delta^{-1/2}] = -\frac{1}{6} R_{\rho\nu}.
\end{equation}
We can obtain an other useful relation taking the contraction and then differentiating Eq.(\ref{diff delta 1}), namely
\begin{equation}
\begin{split}
    d \nabla_\rho \nabla^\nu \nabla_\nu \Delta^{1/2} &= \nabla_\rho \nabla^\nu \nabla_\nu\Delta^{1/2} \sigma_\mu^{\;\,\mu} +\nabla^\nu \nabla_\nu \Delta^{1/2}\sigma_{\mu\;\,\rho}^{\;\,\mu} + \nabla_\rho \nabla_\nu \Delta^{1/2} \sigma_\mu^{\;\,\mu\nu} + \nabla_\nu \Delta^{1/2} \sigma^{\;\,\mu\nu}_{\mu\;\;\,\,\rho}\\
    &+ \nabla_\rho \nabla^\nu \Delta^{1/2} \sigma^{\,\;\mu}_{\mu\;\,\nu} + \nabla^\nu \Delta^{1/2}\sigma_{\mu\,\;\nu\rho}^{\,\;\mu} + \nabla_\rho\Delta^{1/2}\sigma_{\mu\,\;\nu}^{\,\;\mu\,\;\nu}+\Delta^{1/2}\sigma_{\mu\,\;\nu\;\,\rho}^{\,\;\mu\,\;\nu}\\
    &+ 2\sigma_{\;\,\nu\,\;\rho}^{\mu\;\,\nu} \nabla_\mu \Delta^{1/2} +2\sigma^{\mu\;\,\nu}_{\;\,\nu}\nabla_\rho \nabla_\mu \Delta^{1/2} + 2\sigma^\mu_{\;\,\nu\rho}\nabla^\nu \nabla_\mu \Delta^{1/2}\\
    &+ 2 \sigma^\mu_{\;\,\nu}\nabla_\rho \nabla^\nu \nabla_\mu \Delta^{1/2} + 2\sigma^{\mu\nu}_{\;\,\;\,\rho} \nabla_\nu \nabla_\mu \Delta^{1/2} +2\sigma^{\mu\nu} \nabla_\rho \nabla_\nu \nabla_\mu \Delta^{1/2}\\
    &+ 2 \sigma^{\mu}_{\;\,\rho} \nabla^\nu \nabla_\nu \nabla_\mu \Delta^{1/2} + 2 \sigma^\mu \nabla_\rho \nabla^\nu \nabla_\nu \nabla_\mu \Delta^{1/2},
\end{split}
\end{equation}
taking the coincidence limit we get
\begin{equation}\label{contracted 3 delta}
    [\nabla_\rho \nabla^\nu \nabla_\nu \Delta^{1/2}]= -\frac{1}{6}([\sigma_{\mu\,\;\nu\,\;\rho}^{\;\,\mu\,\;\nu}] + [R_\rho^{\,\;\mu}\nabla_\mu \Delta^{1/2}]) = -\frac{1}{6} [\sigma_{\mu\,\;\nu\,\;\rho}^{\;\,\mu\,\;\nu}] = -\frac{1}{6} \nabla_\rho R.
\end{equation}
Another differentiation gives
\begin{equation}
    [\nabla^\nu\nabla_\nu\nabla^\mu\nabla_\mu\Delta^{1/2}] = -\frac{1}{5}\nabla^\mu\nabla_\mu R + \frac{1}{36}R^2- \frac{1}{30}R_{\mu\nu}R^{\mu\nu} + \frac{1}{30}R_{\mu\nu\rho\sigma}R^{\mu\nu\rho\sigma}
\end{equation}
As we did for $\sigma$ and $\Delta$, we can compute coincidence limits for the coefficients $a_k(x,x')$ differentiating and inductively using 
\begin{equation}\label{crucial relation for coefficients}
    ka_k + \sigma^\mu \nabla_\mu a_k + \Delta^{-1/2} \mathcal{O}(\Delta^{1/2} a_{k-1})=0
\end{equation}
with the boundary condition $\sigma^\mu \nabla_\mu a_0$. The recursive equation (\ref{crucial relation for coefficients}) is obtained by (\ref{heat equation}), (\ref{Seely-DeWitt expansion}), together with (\ref{crucial relations}). In the case of a simple scalar field the first coefficient is trivially $a_0(x,x')=1$ because the Seely-De Witt expansion solves the diffusion equation in flat space. For an operator as (\ref{Laplace-type operator}) and for $k=1$ we have
\begin{equation}\label{c.l.coefficient 1}
    [a_1] = [\Delta^{-1/2}\nabla^\mu \partial_\mu \Delta^{1/2} - E] = -\frac{R}{6} -E.
\end{equation}
Not without reason we can continue differentiating (\ref{crucial relation for coefficients}) in the case $k=1$ to also get the coincidence limits for the first and the second derivative of the coefficient $a_1$. The resulting equation for the first derivative is
\begin{equation}
    2[\nabla_\nu a_1] = [\Delta^{-1/2} \nabla_\nu\nabla^\mu\nabla_\mu \Delta^{1/2}]-\nabla_\nu E,
\end{equation}
substituting the result (\ref{contracted 3 delta}), the equation takes the form
\begin{equation}
    [\nabla_\nu a_1] =-\frac{1}{12}\nabla_\nu R - \nabla_\nu E.
\end{equation}
The equation for the second derivative is 
\begin{equation}
    2[\nabla^\nu\nabla_\nu a_1] = [\nabla^\nu\nabla_\nu \Delta^{-1/2}\cdot\nabla^\mu\nabla_\mu \Delta^{1/2}] + [\Delta^{-1/2}\nabla^\nu\nabla_\nu\nabla^\mu\nabla_\mu \Delta^{1/2}] - \nabla^\nu\nabla_\nu E,
\end{equation}
using the previous results for $\Delta$, we get
\begin{equation}\label{seconde derivative c.l. a_1}
    [\nabla^\nu\nabla_\nu a_1] = -\frac{1}{10}\nabla^\mu \nabla_\mu R -\frac{1}{60} R_{\mu\nu}R^{\mu\nu}  + \frac{1}{60} R_{\mu\nu\rho\sigma}R^{\mu\nu\rho\sigma} +\frac{1}{2}\nabla^\mu \nabla_\mu E.
\end{equation}
We will conclude with the computation of the coincidence limit for the second coefficient, in the case $k=2$, (\ref{crucial relation for coefficients}) becomes
\begin{equation}
    2a_2 + \sigma^\mu \nabla_\mu a_2 + \Delta^{-1/2}(-g^{\mu\nu}\nabla_\mu\partial_\nu + E)(\Delta^{1/2}a_1),
\end{equation}
and the coincidence limit is
\begin{equation}
    \begin{split}
        2[a_2] &=[\Delta^{-1/2}g^{\mu\nu} \nabla_\mu \partial_\nu (\Delta^{1/2}a_1) - \Delta^{-1/2}Ea_1]\\
        &= [g^{\mu\nu}(\nabla_\mu \partial_\nu\Delta^{1/2} \cdot a_1 + \partial_\nu \Delta^{1/2}\nabla_\mu a_1 + \nabla_\mu \Delta^{1/2} \partial_\nu a_1 + \Delta^{1/2}\nabla_\mu \partial_\nu a_1) + Ea_1]\\
        &=[\nabla^\mu \partial_\mu \Delta^{1/2} \cdot a_1+\nabla^\mu \partial_\mu a_1 + Ea_1],
    \end{split}
\end{equation}
which using (\ref{c.l.coefficient 1}), (\ref{seconde derivative c.l. a_1}) and (\ref{c.l.delta 2}), becomes
\begin{equation}
    [a_2] = -\frac{1}{72} R^2 + \frac{1}{20} \nabla^\mu \nabla_\mu R + \frac{1}{120} R_{\mu\nu}R^{\mu\nu} - \frac{1}{120}R_{\mu\nu\rho\sigma}R^{\mu\nu\rho\sigma} + \frac{1}{2}E^2 - \frac{1}{2}\nabla^\mu\nabla_\mu E,
\end{equation}
even if it should be something like, ERROR
\begin{equation}
    [a_2] = \frac{1}{72}R^2 + \frac{1}{6}RE + \frac{1}{2} E^2 - \frac{1}{6}\nabla^2
\left(E-\frac{1}{6}R\right)+ \frac{1}{180}(R_{\mu\nu\rho\sigma}R^{\mu\nu\rho\sigma}-R_{\mu\nu} R^{\mu\nu})\end{equation}